\documentclass[11pt,a4paper]{article}
\usepackage[utf8]{inputenc}
\usepackage[T1]{fontenc}
\usepackage[spanish,es-minimal,es-noshorthands]{babel}
\usepackage{lmodern}
\usepackage[sfdefault,lining]{FiraSans}
\usepackage[a4paper,top=2.3cm,bottom=2.2cm,left=1.8cm,right=1.8cm,headheight=1.3cm,headsep=0.4cm,footskip=1cm]{geometry}
\usepackage{amsmath,amssymb}
\usepackage{icomma}
\usepackage{graphicx}
\usepackage[table]{xcolor}
\usepackage{tikz}
\usepackage{fontawesome5}
\usepackage{ccicons}
\usepackage{fancyhdr}
\usepackage{lastpage}
\usepackage{tasks}
\usepackage{needspace}
\usepackage[most]{tcolorbox}
\usepackage{array}
\usepackage[hidelinks]{hyperref}
\AtBeginDocument{\hypersetup{pdftitle={Ficha de ejercicios de la unidad 2: Polinomios y fracciones algebraicas. 4º ESO},pdfsubject={Matemáticas B 4º ESO},pdfkeywords={polinomios, valor numérico, identidades notables, binomio de Newton, Ruffini, teorema del resto, factorización, m.c.d., m.c.m., fracciones algebraicas}}}

\definecolor{navy}{HTML}{1B2748}
\definecolor{accent}{HTML}{E8603A}
\definecolor{teal}{HTML}{1F9E90}
\definecolor{soft}{HTML}{EEF1F8}
\definecolor{muted}{HTML}{5D6685}

\setlength{\parindent}{0pt}
\setlength{\parskip}{3pt}
\settasks{label=\alph*), label-format=\bfseries\color{navy}, label-width=1.3em, item-indent=1.8em, column-sep=0.6em, after-item-skip=5pt, before-skip=3pt, after-skip=4pt}

% ---------- capçalera i peu ----------
\newcommand{\llicencia}{\href{https://creativecommons.org/licenses/by-nc-sa/4.0/deed.es}{\ccbyncsa}\ \ Licencia Creative Commons Atribución-NoComercial-CompartirIgual 4.0 Internacional (CC BY-NC-SA 4.0)}
\pagestyle{fancy}
\fancyhf{}
\fancyhead[L]{\raisebox{-0.35\height}{\includegraphics[height=0.95cm]{logo-gva.png}}}
\fancyhead[C]{\small\color{muted}Ficha de ejercicios · Unidad 2: Polinomios y fracciones algebraicas · 4\textsuperscript{o} ESO B}
\fancyhead[R]{\raisebox{-0.35\height}{\includegraphics[height=0.95cm]{logo-vermellar.png}}}
\fancyfoot[L]{\scriptsize\color{muted}\llicencia}
\fancyfoot[R]{\small\color{muted}\thepage}
\renewcommand{\headrulewidth}{0.4pt}
\renewcommand{\headrule}{\color{muted!40}\hrule height \headrulewidth}
\fancypagestyle{primera}{\fancyhf{}\renewcommand{\headrulewidth}{0pt}
  \fancyfoot[L]{\scriptsize\color{muted}\llicencia}
  \fancyfoot[R]{\small\color{muted}\thepage}}

% ---------- nivells i exercicis ----------
\newcommand{\estrelles}[1]{{\color{accent}\ifcase#1\or\faStar\faStar[regular]\faStar[regular]\or\faStar\faStar\faStar[regular]\or\faStar\faStar\faStar\fi}}
\newcounter{exer}
\newcommand{\ex}[1]{\par\needspace{4\baselineskip}\medskip\refstepcounter{exer}\noindent{\color{navy}\bfseries\theexer.}\ \ {\small\estrelles{#1}}\ \ }

\newcommand{\seccio}[2]{\par\needspace{16\baselineskip}\bigskip
  \begin{tcolorbox}[colback=navy,colframe=navy,coltext=white,arc=2mm,boxrule=0pt,left=3mm,right=3mm,top=1.2mm,bottom=1.2mm]
  {\large\bfseries #1\quad #2}
  \end{tcolorbox}\vspace{-2pt}}

\newcommand{\exf}[1]{\par\medskip\refstepcounter{exer}\noindent{\color{navy}\bfseries\theexer.}\ \ {\small\estrelles{#1}}\ \ }
\newcolumntype{C}[1]{>{\centering\arraybackslash}p{#1}}
% quadre de recordatori (només a la fitxa d'exercicis)
\newcommand{\recorda}[1]{\par\smallskip\begin{tcolorbox}[colback=accent!7,colframe=accent!55,arc=2mm,boxrule=0.5pt,left=2.5mm,right=2.5mm,top=1mm,bottom=1mm]\small{\color{accent}\faLightbulb[regular]}\ \textbf{Recuerda:}\ #1\end{tcolorbox}}
\newcommand{\res}[1]{#1}\newcommand{\rest}[1]{#1}

\def\fmneed{11}% que el títol d'una secció no quede sol al final del full
% =====================================================================
% Capçalera FentMates de les fitxes: un color per a cada curs.
% Els colors són els de la pàgina inicial (LEVELS a assets/js/materials.js);
% el logo ve de _eines/logo/pdf/ (python3 _eines/logo/genera.py).
%
% Ús, al final del preàmbul de la unitat:
%   % =====================================================================
% Capçalera FentMates de les fitxes: un color per a cada curs.
% Els colors són els de la pàgina inicial (LEVELS a assets/js/materials.js);
% el logo ve de _eines/logo/pdf/ (python3 _eines/logo/genera.py).
%
% Ús, al final del preàmbul de la unitat:
%   % =====================================================================
% Capçalera FentMates de les fitxes: un color per a cada curs.
% Els colors són els de la pàgina inicial (LEVELS a assets/js/materials.js);
% el logo ve de _eines/logo/pdf/ (python3 _eines/logo/genera.py).
%
% Ús, al final del preàmbul de la unitat:
%   \input{../comu/fentmates.tex}
%   \fmcurs{1eso}
% i, al començament del document:
%   \capcalera{Matemàtiques}{1r ESO}{Fitxa d'exercicis · Unitat 2}{Divisibilitat}
% Opcional abans de \input: \def\fmneed{14} (espai mínim abans de cada secció).
% =====================================================================
\makeatletter
\def\fm@logodir{../../logo/pdf}
\providecommand{\fmneed}{10}
\def\fm@cursnom{}

% clau del curs, color i color del text damunt del color
\newcommand{\fm@curs}[3]{\@namedef{fm@color@#1}{#2}\@namedef{fm@text@#1}{#3}}
\fm@curs{1eso}{E6B84A}{navy}
\fm@curs{2eso}{3B82F6}{white}
\fm@curs{3eso}{FF7A50}{white}
\fm@curs{4eso-a}{A78BFA}{white}
\fm@curs{4eso-b}{6D28D9}{white}
\fm@curs{1batx-mi}{9B1C1C}{white}
\fm@curs{1batx-ccss}{EF6B6B}{white}
\fm@curs{2batx-mii}{1E6B47}{white}
\fm@curs{2batx-ccss}{2BB3A3}{white}

\newcommand{\fmcurs}[1]{%
  \@ifundefined{fm@color@#1}{\PackageError{fentmates}{Curs desconegut: #1}{}}{}%
  \edef\fm@tmp{\noexpand\definecolor{curs}{HTML}{\@nameuse{fm@color@#1}}%
               \noexpand\colorlet{curstext}{\@nameuse{fm@text@#1}}}\fm@tmp
  \colorlet{cursfosc}{curs!72!black}%
  \def\fm@logo{\fm@logodir/vertical-#1.pdf}%
  \def\fm@marca{\fm@logodir/marca-#1.pdf}}

% ---------- capçalera de la primera pàgina ----------
% #1 matèria (petit), #2 curs, #3 tipus de document i unitat, #4 títol
\newcommand{\capcalera}[4]{%
  \gdef\fm@cursnom{#2}%
  \par\noindent
  \begin{tcolorbox}[colback=curs,colframe=curs,boxrule=0pt,arc=2.2mm,width=3.5cm,height=2.15cm,
      valign=center,left=3mm,right=2mm,top=1mm,bottom=1mm,nobeforeafter,box align=center]
    \color{curstext}{\footnotesize #1}\par\vspace{1pt}{\fontsize{19}{21}\selectfont\bfseries #2\par}%
  \end{tcolorbox}\hspace{5mm}%
  \begin{minipage}[c]{\dimexpr\linewidth-3.5cm-5mm-2.9cm\relax}
    {\color{cursfosc}\small\bfseries\MakeUppercase{#3}}\par\vspace{2pt}
    {\color{navy}\fontsize{20}{23}\selectfont\bfseries #4\par}
  \end{minipage}\hfill
  \begin{minipage}[c]{2.6cm}\raggedleft\includegraphics[width=2.6cm]{\fm@logo}\end{minipage}\par
  \vspace{2pt}}

% ---------- barres de secció del color del curs ----------
\renewcommand{\seccio}[2]{\par\needspace{\fmneed\baselineskip}\bigskip
  \begin{tcolorbox}[colback=curs,colframe=curs,coltext=curstext,arc=2mm,boxrule=0pt,left=3mm,right=3mm,top=1.2mm,bottom=1.2mm]
  {\large\bfseries #1\quad #2}
  \end{tcolorbox}\vspace{-2pt}}

% ---------- estrelles de dificultat del color del curs (a Batxillerat no se'n posen) ----------
\@ifundefined{estrelles}{}{%
  \renewcommand{\estrelles}[1]{{\color{curs}\ifcase#1\or\faStar\faStar[regular]\faStar[regular]\or\faStar\faStar\faStar[regular]\or\faStar\faStar\faStar\fi}}}

% ---------- capçalera i peu de les pàgines ----------
% primera pàgina: logos institucionals al peu
\fancypagestyle{primera}{\fancyhf{}\renewcommand{\headrulewidth}{0pt}
  \fancyfoot[L]{\raisebox{-0.32\height}{\includegraphics[height=0.78cm]{logo-gva.png}}\hspace{5mm}%
    \raisebox{-0.32\height}{\includegraphics[height=0.78cm]{logo-vermellar.png}}}
  \fancyfoot[R]{\scriptsize\color{muted}\href{https://creativecommons.org/licenses/by-nc-sa/4.0/deed.ca}{\ccbyncsa}\ \ CC BY-NC-SA 4.0 · fentmates.org\qquad\small\thepage}}
% la resta: el símbol a l'esquerra i el curs a la dreta, amb el títol al centre com abans
\fancyhead[L]{\raisebox{-0.3\height}{\includegraphics[height=0.62cm]{\fm@marca}}}
\fancyhead[R]{\small\bfseries\color{cursfosc}\fm@cursnom}
\renewcommand{\headrule}{\color{curs!45}\hrule height \headrulewidth}
\makeatother

%   \fmcurs{1eso}
% i, al començament del document:
%   \capcalera{Matemàtiques}{1r ESO}{Fitxa d'exercicis · Unitat 2}{Divisibilitat}
% Opcional abans de \input: \def\fmneed{14} (espai mínim abans de cada secció).
% =====================================================================
\makeatletter
\def\fm@logodir{../../logo/pdf}
\providecommand{\fmneed}{10}
\def\fm@cursnom{}

% clau del curs, color i color del text damunt del color
\newcommand{\fm@curs}[3]{\@namedef{fm@color@#1}{#2}\@namedef{fm@text@#1}{#3}}
\fm@curs{1eso}{E6B84A}{navy}
\fm@curs{2eso}{3B82F6}{white}
\fm@curs{3eso}{FF7A50}{white}
\fm@curs{4eso-a}{A78BFA}{white}
\fm@curs{4eso-b}{6D28D9}{white}
\fm@curs{1batx-mi}{9B1C1C}{white}
\fm@curs{1batx-ccss}{EF6B6B}{white}
\fm@curs{2batx-mii}{1E6B47}{white}
\fm@curs{2batx-ccss}{2BB3A3}{white}

\newcommand{\fmcurs}[1]{%
  \@ifundefined{fm@color@#1}{\PackageError{fentmates}{Curs desconegut: #1}{}}{}%
  \edef\fm@tmp{\noexpand\definecolor{curs}{HTML}{\@nameuse{fm@color@#1}}%
               \noexpand\colorlet{curstext}{\@nameuse{fm@text@#1}}}\fm@tmp
  \colorlet{cursfosc}{curs!72!black}%
  \def\fm@logo{\fm@logodir/vertical-#1.pdf}%
  \def\fm@marca{\fm@logodir/marca-#1.pdf}}

% ---------- capçalera de la primera pàgina ----------
% #1 matèria (petit), #2 curs, #3 tipus de document i unitat, #4 títol
\newcommand{\capcalera}[4]{%
  \gdef\fm@cursnom{#2}%
  \par\noindent
  \begin{tcolorbox}[colback=curs,colframe=curs,boxrule=0pt,arc=2.2mm,width=3.5cm,height=2.15cm,
      valign=center,left=3mm,right=2mm,top=1mm,bottom=1mm,nobeforeafter,box align=center]
    \color{curstext}{\footnotesize #1}\par\vspace{1pt}{\fontsize{19}{21}\selectfont\bfseries #2\par}%
  \end{tcolorbox}\hspace{5mm}%
  \begin{minipage}[c]{\dimexpr\linewidth-3.5cm-5mm-2.9cm\relax}
    {\color{cursfosc}\small\bfseries\MakeUppercase{#3}}\par\vspace{2pt}
    {\color{navy}\fontsize{20}{23}\selectfont\bfseries #4\par}
  \end{minipage}\hfill
  \begin{minipage}[c]{2.6cm}\raggedleft\includegraphics[width=2.6cm]{\fm@logo}\end{minipage}\par
  \vspace{2pt}}

% ---------- barres de secció del color del curs ----------
\renewcommand{\seccio}[2]{\par\needspace{\fmneed\baselineskip}\bigskip
  \begin{tcolorbox}[colback=curs,colframe=curs,coltext=curstext,arc=2mm,boxrule=0pt,left=3mm,right=3mm,top=1.2mm,bottom=1.2mm]
  {\large\bfseries #1\quad #2}
  \end{tcolorbox}\vspace{-2pt}}

% ---------- estrelles de dificultat del color del curs (a Batxillerat no se'n posen) ----------
\@ifundefined{estrelles}{}{%
  \renewcommand{\estrelles}[1]{{\color{curs}\ifcase#1\or\faStar\faStar[regular]\faStar[regular]\or\faStar\faStar\faStar[regular]\or\faStar\faStar\faStar\fi}}}

% ---------- capçalera i peu de les pàgines ----------
% primera pàgina: logos institucionals al peu
\fancypagestyle{primera}{\fancyhf{}\renewcommand{\headrulewidth}{0pt}
  \fancyfoot[L]{\raisebox{-0.32\height}{\includegraphics[height=0.78cm]{logo-gva.png}}\hspace{5mm}%
    \raisebox{-0.32\height}{\includegraphics[height=0.78cm]{logo-vermellar.png}}}
  \fancyfoot[R]{\scriptsize\color{muted}\href{https://creativecommons.org/licenses/by-nc-sa/4.0/deed.ca}{\ccbyncsa}\ \ CC BY-NC-SA 4.0 · fentmates.org\qquad\small\thepage}}
% la resta: el símbol a l'esquerra i el curs a la dreta, amb el títol al centre com abans
\fancyhead[L]{\raisebox{-0.3\height}{\includegraphics[height=0.62cm]{\fm@marca}}}
\fancyhead[R]{\small\bfseries\color{cursfosc}\fm@cursnom}
\renewcommand{\headrule}{\color{curs!45}\hrule height \headrulewidth}
\makeatother

%   \fmcurs{1eso}
% i, al començament del document:
%   \capcalera{Matemàtiques}{1r ESO}{Fitxa d'exercicis · Unitat 2}{Divisibilitat}
% Opcional abans de \input: \def\fmneed{14} (espai mínim abans de cada secció).
% =====================================================================
\makeatletter
\def\fm@logodir{../../logo/pdf}
\providecommand{\fmneed}{10}
\def\fm@cursnom{}

% clau del curs, color i color del text damunt del color
\newcommand{\fm@curs}[3]{\@namedef{fm@color@#1}{#2}\@namedef{fm@text@#1}{#3}}
\fm@curs{1eso}{E6B84A}{navy}
\fm@curs{2eso}{3B82F6}{white}
\fm@curs{3eso}{FF7A50}{white}
\fm@curs{4eso-a}{A78BFA}{white}
\fm@curs{4eso-b}{6D28D9}{white}
\fm@curs{1batx-mi}{9B1C1C}{white}
\fm@curs{1batx-ccss}{EF6B6B}{white}
\fm@curs{2batx-mii}{1E6B47}{white}
\fm@curs{2batx-ccss}{2BB3A3}{white}

\newcommand{\fmcurs}[1]{%
  \@ifundefined{fm@color@#1}{\PackageError{fentmates}{Curs desconegut: #1}{}}{}%
  \edef\fm@tmp{\noexpand\definecolor{curs}{HTML}{\@nameuse{fm@color@#1}}%
               \noexpand\colorlet{curstext}{\@nameuse{fm@text@#1}}}\fm@tmp
  \colorlet{cursfosc}{curs!72!black}%
  \def\fm@logo{\fm@logodir/vertical-#1.pdf}%
  \def\fm@marca{\fm@logodir/marca-#1.pdf}}

% ---------- capçalera de la primera pàgina ----------
% #1 matèria (petit), #2 curs, #3 tipus de document i unitat, #4 títol
\newcommand{\capcalera}[4]{%
  \gdef\fm@cursnom{#2}%
  \par\noindent
  \begin{tcolorbox}[colback=curs,colframe=curs,boxrule=0pt,arc=2.2mm,width=3.5cm,height=2.15cm,
      valign=center,left=3mm,right=2mm,top=1mm,bottom=1mm,nobeforeafter,box align=center]
    \color{curstext}{\footnotesize #1}\par\vspace{1pt}{\fontsize{19}{21}\selectfont\bfseries #2\par}%
  \end{tcolorbox}\hspace{5mm}%
  \begin{minipage}[c]{\dimexpr\linewidth-3.5cm-5mm-2.9cm\relax}
    {\color{cursfosc}\small\bfseries\MakeUppercase{#3}}\par\vspace{2pt}
    {\color{navy}\fontsize{20}{23}\selectfont\bfseries #4\par}
  \end{minipage}\hfill
  \begin{minipage}[c]{2.6cm}\raggedleft\includegraphics[width=2.6cm]{\fm@logo}\end{minipage}\par
  \vspace{2pt}}

% ---------- barres de secció del color del curs ----------
\renewcommand{\seccio}[2]{\par\needspace{\fmneed\baselineskip}\bigskip
  \begin{tcolorbox}[colback=curs,colframe=curs,coltext=curstext,arc=2mm,boxrule=0pt,left=3mm,right=3mm,top=1.2mm,bottom=1.2mm]
  {\large\bfseries #1\quad #2}
  \end{tcolorbox}\vspace{-2pt}}

% ---------- estrelles de dificultat del color del curs (a Batxillerat no se'n posen) ----------
\@ifundefined{estrelles}{}{%
  \renewcommand{\estrelles}[1]{{\color{curs}\ifcase#1\or\faStar\faStar[regular]\faStar[regular]\or\faStar\faStar\faStar[regular]\or\faStar\faStar\faStar\fi}}}

% ---------- capçalera i peu de les pàgines ----------
% primera pàgina: logos institucionals al peu
\fancypagestyle{primera}{\fancyhf{}\renewcommand{\headrulewidth}{0pt}
  \fancyfoot[L]{\raisebox{-0.32\height}{\includegraphics[height=0.78cm]{logo-gva.png}}\hspace{5mm}%
    \raisebox{-0.32\height}{\includegraphics[height=0.78cm]{logo-vermellar.png}}}
  \fancyfoot[R]{\scriptsize\color{muted}\href{https://creativecommons.org/licenses/by-nc-sa/4.0/deed.ca}{\ccbyncsa}\ \ CC BY-NC-SA 4.0 · fentmates.org\qquad\small\thepage}}
% la resta: el símbol a l'esquerra i el curs a la dreta, amb el títol al centre com abans
\fancyhead[L]{\raisebox{-0.3\height}{\includegraphics[height=0.62cm]{\fm@marca}}}
\fancyhead[R]{\small\bfseries\color{cursfosc}\fm@cursnom}
\renewcommand{\headrule}{\color{curs!45}\hrule height \headrulewidth}
\makeatother

\fmcurs{4eso-b}

\begin{document}
\thispagestyle{primera}

\capcalera{Matemáticas B}{4\textsuperscript{o} ESO}{Ficha de ejercicios · Unidad 2}{\raggedright Polinomios y fracciones algebraicas}
\vspace{10pt}
Nombre y apellidos: \rule{6.6cm}{0.4pt}\hfill Grupo: \rule{1.6cm}{0.4pt}\hfill Fecha: \rule{2.2cm}{0.4pt}
\vspace{8pt}

\begin{tcolorbox}[colback=soft,colframe=soft,arc=2mm,left=3mm,right=3mm,top=1.5mm,bottom=1.5mm]
\small Ejercicios graduados por niveles de dificultad:\quad \estrelles{1}\ básico \qquad \estrelles{2}\ medio \qquad \estrelles{3}\ avanzado.\\[2pt] Resuélvelos en el cuaderno mostrando todos los pasos. Ordena siempre los polinomios y, en las divisiones, deja el lugar de los términos que faltan. Factoriza antes de operar con fracciones y simplifica todos los resultados. Recoge los ejercicios de las fichas del departamento de la unidad.
\end{tcolorbox}

% =====================================================
\seccio{1}{Polinomios y valor numérico}

\exf{1} Completa la tabla siguiente:\begin{center}\renewcommand{\arraystretch}{1.45}\small\begin{tabular}{|l|C{1cm}|C{1.9cm}|C{1.5cm}|C{2cm}|C{1.7cm}|}\hline\rowcolor{soft}\textbf{Polinomio} & \textbf{Grado} & \textbf{Térm. indep.} & \textbf{N.\textsuperscript{o} términos} & \textbf{¿Completo?} & \textbf{Coef. princ.}\\\hline$P(x)=5x^4-3x^3+2x+10$ &  &  &  &  & \\\hline$Q(x)=-5x^5-\frac32x^4-x^3+9x^2-6$ &  &  &  &  & \\\hline$R(x)=4x^2+x+1$ &  &  &  &  & \\\hline$S(x)=5x^2+2x^5+9x$ &  &  &  &  & \\\hline$T(x)=3x-2$ &  &  &  &  & \\\hline$U(x)=x^3-4x^2+x+1$ &  &  &  &  & \\\hline\end{tabular}\end{center}

\ex{1} Dado el polinomio $P(x)=-3x^4-2x^3-x^2-1$, calcula su valor numérico para $x=2$ y $x=-2$.

\ex{1} Calcula el valor numérico de los polinomios siguientes para $x=3$, $x=-1$ y $x=0$:\begin{tasks}(2)
  \task $P(x)=6x^3-x^2-3$
  \task $Q(x)=-3x^4+2x^3+7x^2-x+5$
  \task $R(x)=8$
  \task $S(x)=x^3-1$
\end{tasks}

\ex{2} ¿Qué valor debe tener $k$ para que el valor numérico del polinomio $P(x)=x^4-2x^2+5x+k$ para $x=2$ sea 6?

\ex{2} ¿Cuánto vale $a$ en el polinomio $P(x)=-x^4-4x^3+ax+10$, si sabemos que el valor numérico para $x=-3$ es 10?


% =====================================================
\seccio{2}{Operaciones e identidades notables}

\exf{2} Sean $P(x)=x^3+4x^2-3x+1$, $Q(x)=2x^3-x+3$ y $R(x)=4x^2-x$. Calcula:\begin{tasks}(2)
  \task $[P(x)-3\cdot Q(x)]\cdot R(x)$
  \task $[P(x)-R(x)]\cdot Q(x)$
\end{tasks}

\ex{1} Efectúa las operaciones siguientes aplicando las identidades notables:\begin{tasks}(3)
  \task $(x+3)^2$
  \task $(5x^2+7)^2$
  \task $(2x-3)^2$
  \task $(2x-5)\cdot(2x+5)$
  \task $\left(\sqrt2x-\sqrt{50}\right)^2$
  \task $\left(\frac23x-\frac32\right)^2$
\end{tasks}\recorda{$(a\pm b)^2=a^2\pm2ab+b^2$ \quad y \quad $(a+b)(a-b)=a^2-b^2$.}

\ex{2} Calcula:\begin{tasks}(1)
  \task $3x(2x-1)-(x-3)(x+3)+(x-2)^2$
  \task $(2x-1)^2+(x-1)(3-x)-3(x+5)^2$
  \task $\frac43(x-3)^2-\frac13(3x-1)(3x+1)-\frac13(4x^3+35)$
\end{tasks}


% =====================================================
\seccio{3}{Binomio de Newton}

\exf{2} Calcula las potencias siguientes utilizando la fórmula del binomio de Newton:\begin{tasks}(3)
  \task $(x+2)^3$
  \task $(x-1)^4$
  \task $(2x+3)^3$
  \task $(2x+1)^5$
  \task $(2x-3)^3$
  \task $(-2+x)^6$
\end{tasks}\recorda{$(a+b)^n=\binom n0a^n+\binom n1a^{n-1}b+\cdots+\binom nnb^n$; los coeficientes son una fila del triángulo de Tartaglia. Si hay una resta, los signos se alternan.}

\ex{3} Calcula las potencias siguientes:\begin{tasks}(4)
  \task $(x^2+x)^3$
  \task $(2x^4-y)^5$
  \task $(4x^3+2)^4$
  \task $(x^6-3y^2)^5$
  \task $(3y^2+5x)^3$
  \task $(-2x^3+2x)^4$
  \task $(-2x^3-3x)^6$
  \task $(x-2y)^4$
\end{tasks}


% =====================================================
\seccio{4}{División de polinomios y regla de Ruffini}

\exf{2} Efectúa las divisiones de polinomios siguientes:\begin{tasks}(1)
  \task $(3x^4+2x^3-7x^2+5x+12):(x^2-2x+1)$
  \task $(x^5-3x^4+x^2+3x-5):(x^2+x+1)$
  \task $(4x^5+3x^3+7x^2+6x+8):(x^2+1)$
  \task $(x^5-2x^4+3x^3-5x^2+10x-15):(x^2-2x+3)$
  \task $(18x^6-3x^4+2x^3+x):(x^2-2)$
\end{tasks}\recorda{Deja vacío el lugar de los términos que faltan en el dividendo. Comprobación: $D(x)=d(x)\cdot c(x)+r(x)$.}

\ex{1} Haz las divisiones con la regla de Ruffini e indica el cociente y el resto:\begin{tasks}(1)
  \task $(x^3-6x^2+11x-5):(x-3)$
  \task $(x^4+7x^3+13x^2-x-17):(x+4)$
  \task $(x^7-9x^6+19x^5+12x^4-3x^3+19x^2-37x-37):(x-5)$
\end{tasks}

\ex{1} Determina el resto de las divisiones siguientes con la regla de Ruffini:\begin{tasks}(3)
  \task $(2x^3-5x^2+3):(x-2)$
  \task $(x^4-3x^3+1):(x+2)$
  \task $(x^4+1):(x+1)$
\end{tasks}


% =====================================================
\seccio{5}{Valor numérico y teorema del resto}

\exf{2} Calcula $m$ para que el polinomio $P(x)=x^3-mx^2+5x-2$ sea divisible por $(x+1)$.

\ex{2} El resto de la división $(2x^4+kx^3-7x+6):(x-2)$ es igual a $-8$. ¿Cuánto vale $k$?

\ex{1} Calcula el valor numérico de $P(x)=3x^4-x^2+2x-5$ para $x=2$. Hazlo de dos formas diferentes.

\ex{2} Mediante el teorema del resto, comprueba si es verdad o no que:\begin{tasks}(2)
  \task $(x-1)$ es divisor de $x^{55}+1$
  \task $(x+1)$ es divisor de $x^{55}+1$
\end{tasks}

\ex{2} Mediante el teorema del resto, comprueba si $x=3$ es solución de la ecuación $6x^5-19x^4+4x^3-8x^2+17x-6=0$.

\ex{2} Mediante el teorema del resto, comprueba si $x=-4$ es solución de la ecuación $2x^6+12x^5+11x^4-25x^3-19x^2-2x-23=0$.


% =====================================================
\seccio{6}{Raíces de un polinomio}

\exf{2} Halla las raíces enteras de los polinomios siguientes:\begin{tasks}(2)
  \task $x^3-6x^2+11x-6$
  \task $x^3+2x^2-x-2$
  \task $x^4-10x^2+9$
  \task $x^3-3x^2-4x+12$
\end{tasks}\recorda{Las raíces enteras de un polinomio con coeficientes enteros son divisores del término independiente.}

\ex{2} Escribe los candidatos a raíz entera de $x^3+x^2-12$ y halla cuáles lo son.

\ex{2} Escribe un polinomio de grado 3 que tenga por raíces $1$, $-2$ y $3$, y uno de grado 2 con una raíz doble en $x=4$.


% =====================================================
\seccio{7}{Factorización}

\exf{1} Saca factor común:\begin{tasks}(2)
  \task $5xy^2+10y$
  \task $12x^3-4x^2$
  \task $6b^2-24$
  \task $4ab^2c^4-28a^2b^3c^5+42a^5b^6c^3$
  \task $24z^2-12z+12$
  \task $5w^4t^6+25w^6t^4$
\end{tasks}

\ex{2} Factoriza los polinomios siguientes:\begin{tasks}(2)
  \task $x^4-x^3-7x^2+x+6$
  \task $x^5+x^4-5x^3-x^2+8x-4$
  \task $x^4+4x^3-2x^2-12x+9$
  \task $2x^5-6x^4-19x^3+24x^2+32x$
  \task $x^3+3x^2+5x+6$
  \task $2x^4-7x^3+x^2+7x-3$
\end{tasks}\recorda{Saca factor común, busca raíces entre los divisores del término independiente con Ruffini y, cuando llegues a grado 2, resuelve la ecuación.}

\ex{2} Halla la factorización de los polinomios siguientes. Si algún factor de grado 2 no tiene raíces reales, es irreducible:\begin{tasks}(2)
  \task $x^3+2x^2+2x+1$
  \task $x^3+4x^2+6x+4$
  \task $x^3+x-2$
  \task $x^3-x-6$
  \task $x^4-29x^2+100$
  \task $x^4-24x^2-25$
  \task $x^4-3x^3+4x^2+3x-5$
  \task $x^4+2x^3+x^2-8x-20$
\end{tasks}


% =====================================================
\seccio{8}{Divisibilidad. M.c.d. y m.c.m.}

\exf{1} Indica, sin hacer la división, si el primer polinomio es divisor del segundo:\begin{tasks}(2)
  \task $x-1$ y $x^3-1$
  \task $x+2$ y $x^2+4$
  \task $x+1$ y $x^4-1$
  \task $x-3$ y $x^2-9$
\end{tasks}

\ex{2} Calcula el m.c.d. y el m.c.m. de:\begin{tasks}(1)
  \task $x^2-4$ y $x^2+4x+4$
  \task $x^2-x$ y $x^2-1$
  \task $x^3-x$ y $x^2-2x+1$
  \task $x^3-2x^2-x+2$ y $x^2-3x+2$
\end{tasks}\recorda{Factoriza primero. M.c.d.: factores comunes con el menor exponente. M.c.m.: todos los factores con el mayor exponente.}


% =====================================================
\seccio{9}{Simplificación de fracciones algebraicas}

\exf{2} Simplifica las fracciones algebraicas siguientes:\begin{tasks}(4)
  \task $\dfrac{x^3+5x^2+8x+4}{x^2-4}$
  \task $\dfrac{x^3-1}{x^2-1}$
  \task $\dfrac{x^3-5x^2-9x+45}{x^3+x^2-21x-45}$
  \task $\dfrac{x^2+9}{x^4-81}$
  \task $\dfrac{x^5+x^4}{x^4-x^2}$
  \task $\dfrac{x^3+x}{x^4-x}$
  \task $\dfrac{x^2-14x+49}{x^2-49}$
  \task $\dfrac{x^3+3x}{x^2+3}$
\end{tasks}\recorda{Factoriza el numerador y el denominador (factor común, identidades notables, Ruffini) y elimina los factores comunes.}

\ex{2} Simplifica:\begin{tasks}(3)
  \task $\dfrac{2x^3-5x^2-3x}{x^2-9}$
  \task $\dfrac{x^3-8}{x^2+2x+4}$
  \task $\dfrac{6x^2-7x-20}{2x^2-18}$
  \task $\dfrac{x^4-16}{x^3+4x^2-4x-16}$
  \task $\dfrac{x^4-1}{x^4-2x^3+2x^2-2x+1}$
\end{tasks}


% =====================================================
\seccio{10}{Multiplicación y división de fracciones algebraicas}

\exf{2} Haz las multiplicaciones y divisiones siguientes y simplifica el resultado:\begin{tasks}(2)
  \task $\dfrac{x-6}{x^2+3x+2}\cdot\dfrac{x^2+2x}{x^2-36}$
  \task $\dfrac{x^2-3x+2}{x^4-9x^2}\cdot\dfrac{x^2-3x}{x-1}$
  \task $\dfrac{2x^2+8x}{x^2-10x+25}\cdot\dfrac{x^2-25}{2x}$
  \task $\dfrac{\ \frac{1}{x-1}\ }{\frac{2x}{x+2}}\cdot\dfrac{\ \frac{x-1}{x+2}\ }{\frac1x}$
  \task $\dfrac{x^3-5x^2}{x^3+3x^2-4x-12}:\dfrac{x-5}{x^3+4x^2+4x}$
  \task $\dfrac{x^2-3x-4}{x^3+x^2-x-1}:\dfrac{x^3-2x^2-7x-4}{x^2-1}$
  \task $\dfrac{x^3}{8x-8}:\dfrac{x^2}{x^2-x}$
  \task $\dfrac{\ \frac{x}{x-3}\ }{\frac{x+3}{x}}:\dfrac{\ \frac{1}{x+1}\ }{\frac{x^2-9}{x^3+x^2}}$
  \task $\dfrac{x^2-6x+5}{x^2+5x+6}\cdot\dfrac{2x^2-8}{x^2-x}:\dfrac{2x-10}{x^2+3x}$
\end{tasks}\recorda{Factoriza antes de multiplicar: así verás qué factores se simplifican.}


% =====================================================
\seccio{11}{Suma y resta de fracciones algebraicas}

\exf{2} Efectúa las sumas y restas siguientes de fracciones algebraicas:\begin{tasks}(2)
  \task $\dfrac{5x+1}{x-2}+\dfrac{x^2+x-3}{x-2}$
  \task $\dfrac{x^2-1}{3x+2}-\dfrac{3x^2-5x-10}{3x+2}$
  \task $\dfrac{1}{x^2-1}-\dfrac{2}{x^2-2x-3}$
  \task $\dfrac{1}{x-5}+\dfrac{1}{x+5}$
  \task $\dfrac{x-3}{x}+\dfrac{6}{x^2+2x}-\dfrac{x+1}{x+2}$
  \task $\dfrac{x+1}{x^3-1}-\dfrac{x-2}{x^2-1}$
  \task $\dfrac{x+1}{x-1}+\dfrac{x-1}{x+1}$
  \task $\dfrac{2x}{x^2-16}+\dfrac{x+1}{3x+12}$
  \task $\dfrac2x+\dfrac{5x}{x-3}-\dfrac{x-4}{x^2-9}$
  \task $\dfrac{x-2}{x^3-x}+\dfrac{x}{x^2-1}$
\end{tasks}\recorda{Factoriza los denominadores, calcula su m.c.m. y reduce a común denominador. Simplifica el resultado si se puede.}

\ex{3} Opera y simplifica:\begin{tasks}(2)
  \task $\dfrac{2}{x-2}-\dfrac{6}{x^2-x-2}+\dfrac{1}{x+1}$
  \task $\dfrac{2x^2-x}{x+3}+\dfrac{2x}{x-3}+\dfrac{12}{x^2-9}$
  \task $\dfrac{x}{x+1}-\dfrac{16}{6}+\dfrac{x+1}{x}$
  \task $\dfrac{x+1}{x^2}-\dfrac{x}{x+1}-\dfrac{1}{x^3+x^2}$
\end{tasks}


% =====================================================
\seccio{12}{Operaciones combinadas}

\exf{3} Opera y simplifica:\begin{tasks}(1)
  \task $\dfrac{a+b}{a-b}-\dfrac{a-b}{a+b}+\dfrac{a^2}{a^2-b^2}$
  \task $\left(1+\dfrac1x\right):\left(1+\dfrac1{x+1}\right):\dfrac{x+1}{x+2}$
  \task $(a^2-b^2):\left(\dfrac1a+\dfrac1b\right)$
  \task $\left(x-\dfrac{x^2}{x-y}\right):\left(y-\dfrac{y^2}{x+y}\right)$
  \task $\dfrac{3x+5}{x^2-4}+\dfrac{2}{2-x}+\dfrac{-x^2+4x-1}{x-2}:(x+2)$
\end{tasks}\recorda{Respeta la jerarquía: primero los paréntesis, después los productos y cocientes (de izquierda a derecha) y al final las sumas y restas.}


% =====================================================
\seccio{13}{Problemas}

\exf{2} Un rectángulo tiene una base de $x+3$ cm y un área de $x^2+5x+6$ cm$^2$. Expresa su altura y su perímetro.

\ex{2} El volumen de una caja es $x^3+6x^2+8x$ cm$^3$ y la base es un rectángulo de lados $x$ y $x+2$. ¿Cuál es la altura?

\vfill
\begin{center}\color{muted}\small\faCheckCircle[regular]\ \ ¡Ánimo! Comprueba tus resultados con las herramientas interactivas de los apuntes de la unidad.\end{center}

\end{document}
